\subsection{Discrete Representation Spaces}

In discrete representation spaces, information is represented through a finite set of distinct states:

\begin{equation}
    \mathcal{Z}_{D} = \{z_1, z_2, ..., z_n\}
\end{equation}

where each representation $z_i$ corresponds to a specific symbol, category, or pattern. Relationships between different states are not inherently represented, meaning that two states are either identical or distinct.

This type of representation is commonly used in symbolic systems and early associative memory models. For example, classical Hopfield networks represent memories as discrete binary patterns:

\begin{equation}
    \xi_i \in \{-1,+1\}^{N}
\end{equation}

where each stored pattern corresponds to a stable state of the network~\cite{hopfield1982}.

While discrete representations allow for clear separation between individual states, they provide limited capability for expressing gradual relationships or similarities between different representations.
